\documentclass[conference]{IEEEtran}
\usepackage[utf8]{inputenc}
\usepackage[T1]{fontenc}
\usepackage[portuguese]{babel}
\usepackage{cite}
\usepackage{amsmath,amssymb,amsfonts}
\usepackage{algorithmic}
\usepackage{graphicx}
\usepackage{textcomp}
\usepackage{xcolor}
\usepackage{booktabs}
\usepackage{float}
\usepackage{listings}
\usepackage{subcaption}
\usepackage{url}
\usepackage[hidelinks]{hyperref}
\raggedbottom

% Configuração para snippets de código MATLAB/Octave
\lstset{
  language=Matlab,
  basicstyle=\ttfamily\footnotesize,
  keywordstyle=\color{blue},
  commentstyle=\color{gray},
  numbers=left,
  numberstyle=\tiny,
  frame=single,
  breaklines=true
}

\begin{document}

\title{Estudo Comparativo de Arquiteturas de Redes Neuronais e Algoritmos de Otimização em Tarefas de Regressão}

\author{\IEEEauthorblockN{Guilherme Basílio - 2100123}
\IEEEauthorblockA{\textit{Faculdade de Ciências Exatas e da Engenharia} \\
\textit{Universidade da Madeira}\\
Funchal, Portugal \\
2100123@student.uma.pt}
\and
\IEEEauthorblockN{Carlos Santos - 2128623}
\IEEEauthorblockA{\textit{Faculdade de Ciências Exatas e da Engenharia} \\
\textit{Universidade da Madeira}\\
Funchal, Portugal \\
2128623@student.uma.pt}
}

\maketitle

\begin{abstract}

\end{abstract}

\begin{IEEEkeywords}

\end{IEEEkeywords}

\section{Introdução}


\section{Metodologia}

\subsection{Aquisição de dados e inspeção inicial}
\subsection{Análise Exploratória dos Dados}
\subsection{Pré-processamento e Feature Engineering}
\subsection{Particionamento dos Dados}
\subsection{Implementação e Verificação}
\subsection{Ajuste de hiperparamêtros}
\subsection{Experimentos Finais}

\section{Resultados}

\section{Discussão}

\section{Conclusão}


\begin{thebibliography}{00}

\bibitem{b1} F. R. S. Mendonça e F. M. R. M. F. Dias, ``Artificial Neural Networks -- Assignment 2 (2026/27): Regression task: network architectures and training algorithms,'' Faculdade de Ciências Exatas e da Engenharia, Universidade da Madeira, Funchal, Portugal, 2026.

\bibitem{b2} I.-C. Yeh, ``Concrete Compressive Strength,'' UCI Machine Learning Repository, 1998. \url{https://archive.ics.uci.edu/dataset/165/concrete+compressive+strength}

\bibitem{b3} I.-C. Yeh, ``Modeling of strength of high-performance concrete using artificial neural networks,'' \textit{Cement and Concrete Research}, vol. 28, no. 12, pp. 1797--1808, Dec. 1998.

\bibitem{b4} P. Brabban, ``Daily Minimum Temperatures in Melbourne,'' Kaggle Dataset. \url{https://www.kaggle.com/datasets/paulbrabban/daily-minimum-temperatures-in-melbourne}

\bibitem{b5} D. P. Kingma e J. Ba, ``Adam: A Method for Stochastic Optimization,'' in \textit{Proc. 3rd Int. Conf. Learn. Representations (ICLR)}, San Diego, CA, USA, 2015.

\end{thebibliography}

\end{document}